\begin{trapbox}

	\begin{description}[style=nextline, leftmargin=2.8em, labelsep=0.5em, itemsep=0.35em]
		\item[\textbf{\textcolor{CTrap}{易错点一（忽视单调性）}}]
			只验证 $f(m)f(n)<0$，没有检查函数是否严格单调，就判定存在唯一根。
		\item[\textbf{\textcolor{CTrap}{正确理解（需验证单调）}}]
			唯一性依赖于严格单调性，因此除连续和异号外，还必须确认函数在区间上严格单调。
		\item[\textbf{\textcolor{CTrap}{错因分析（缺少唯一条件）}}]
			零点定理只保证存在性，唯一性需由单调性提供。缺单调可能产生多个根或唯一性无法保证。
	\end{description}

	\medskip
	\begin{description}[style=nextline, leftmargin=2.8em, labelsep=0.5em, itemsep=0.35em]
		\item[\textbf{\textcolor{CTrap}{易错点二（混淆根区间）}}]
			误认为根在闭区间 $[m,n]$ 内，或认为端点可能是根；在整点问题中认为根可能在整数点 $k$ 或 $k+1$ 上。
		\item[\textbf{\textcolor{CTrap}{正确理解（根在开区间）}}]
			结论明确指出根在开区间 $(m,n)$ 内，不包括端点；整数区间推论得到根在 $(k,k+1)$ 内，故必不是整数。
		\item[\textbf{\textcolor{CTrap}{错因分析（忽略端点条件）}}]
			充要条件要求 $f(m)f(n)<0$，端点为零时乘积为0，不满足条件，此时根可能在端点，但不在开区间内。
	\end{description}

	\medskip
	\begin{description}[style=nextline, leftmargin=2.8em, labelsep=0.5em, itemsep=0.35em]
		\item[\textbf{\textcolor{CTrap}{易错点三（忽略连续性）}}]
			只知道函数严格单调且端点异号，但未检查连续性，就断言有唯一根。
		\item[\textbf{\textcolor{CTrap}{正确理解（连续前提不可缺）}}]
			零点定理要求函数连续；若函数在区间内有间断点，即使单调且端点异号，也可能没有根。
		\item[\textbf{\textcolor{CTrap}{错因分析（缺乏连续保证）}}]
			不连续时函数可能发生跳变，导致端点符号改变但实际并未经过 $x$ 轴，故无法保证根的存在。
	\end{description}

\end{trapbox}
